% 高光谱智能解译 实验报告 LaTeX 骨架
% 用法：pandoc 实验报告-终稿.md -o 实验报告-终稿.docx  （或直接用本 tex 编译）
% 本骨架仅含结构与关键公式，正文与表格请以 Markdown 终稿为准。

\documentclass[12pt,a4paper]{article}
\usepackage[UTF8]{ctex}
\usepackage{amsmath,amssymb}
\usepackage{booktabs}
\usepackage{graphicx}
\usepackage{geometry}
\geometry{left=2.5cm,right=2.5cm,top=2.5cm,bottom=2.5cm}
\graphicspath{{figures/}}

\title{本科实验报告：大作业5 高光谱智能解译}
\date{}

\begin{document}

\section*{封面（本科实验报告）}
实验名称：大作业5：高光谱智能解译\\
课程名称：人工智能基础理论与技术（人工智能导论）\\
实验时间：课后 \quad 任课教师：刘欢 \quad 实验地点：【请填写】 \quad 助教：阮子航\\
实验类型：$\square$ 原理验证 \quad $\blacksquare$ 综合设计 \quad $\square$ 自主创新\\
姓名（组长）/学号/学院/组员/组号/成绩：【请填写】

\section{实验目的}
\subsection{（一）基本任务}
1. 高光谱数据读取与预处理；2. 高光谱数据降维处理（PCA/LDA）；3. HybridSN 高光谱图像分类。
\subsection{（二）进阶任务}
1. 改进 HybridSN 分类模型；2. 提升高光谱分类精度并与原始模型比较。
\subsection{（三）提高任务}
1. 复现其他高光谱分类论文；2. 多模型对比实验；3. 多数据集验证；4. 进一步模型创新（选做）。

\section{实验环境}
1. 数据（Pavia University / Indian Pines / Salinas）；2. 软件及深度学习环境（Python 3.12、PyTorch 2.12.1+cu130、NumPy/SciPy、Scikit-learn、Matplotlib、RTX 5070）；3. 实验参考程序；4. 实验评价（OA/AA/Kappa）。

\section{实验原理}
\subsection{标准化与降维}
\begin{align}
z_b &= \frac{x_b - \mu_b}{\sigma_b} \\
\Sigma &= \frac{1}{N-1}\sum_{i=1}^{N}(\mathbf{z}_i - \bar{\mathbf{z}})(\mathbf{z}_i - \bar{\mathbf{z}})^{\top},\quad
\Sigma \mathbf{w}_k = \lambda_k \mathbf{w}_k \\
\mathbf{S}_b &= \sum_{c=1}^{C} N_c (\boldsymbol{\mu}_c - \boldsymbol{\mu})(\boldsymbol{\mu}_c - \boldsymbol{\mu})^{\top},\quad
\mathbf{S}_w = \sum_{c=1}^{C} \sum_{i \in c}(\mathbf{x}_i - \boldsymbol{\mu}_c)(\mathbf{x}_i - \boldsymbol{\mu}_c)^{\top}
\end{align}

\subsection{卷积与激活}
\begin{equation}
y_{j}^{l}(x,y,z) = f\Big( \sum_{m} \sum_{p,q,r} w_{j,m}^{l}(p,q,r)\, x_{m}^{l-1}(x+p,\, y+q,\, z+r) + b_j^l \Big),\quad
f(x)=\max(0,x)
\end{equation}

\subsection{Softmax 与交叉熵}
\begin{equation}
P(y=c \mid \mathbf{x}) = \frac{\exp(z_c)}{\sum_{k=1}^{C}\exp(z_k)},\qquad
\mathcal{L} = -\sum_{c=1}^{C} y_c \log P(y=c \mid \mathbf{x})
\end{equation}

\subsection{BatchNorm、残差与全局平均池化}
\begin{align}
\hat{x}^{(k)} &= \frac{x^{(k)} - \mu^{(k)}}{\sqrt{(\sigma^{(k)})^2 + \epsilon}},\quad y^{(k)}=\gamma^{(k)}\hat{x}^{(k)}+\beta^{(k)} \\
\mathbf{y} &= \mathrm{ReLU}\big(\mathrm{BN}(\mathrm{Conv}(\mathbf{x})) + \mathrm{shortcut}(\mathbf{x})\big) \\
g_c &= \frac{1}{HW}\sum_{h=1}^{H}\sum_{w=1}^{W} f_c(h,w)
\end{align}

\subsection{评价指标}
\begin{equation}
\mathrm{OA}=\frac{1}{N}\sum_i \mathbb{1}[\hat{y}_i=y_i],\quad
\mathrm{AA}=\frac{1}{C}\sum_{c=1}^{C}\mathrm{acc}_c,\quad
\kappa=\frac{p_o-p_e}{1-p_e}
\end{equation}

\section{实验内容和结果分析}
（实验过程、代码截图、实验结果、误差分析、图文并茂，另提交表格/代码/结果/参考文献作为附件。数值以正式运行产出的 metrics.json / summary_row.json 为准。）

\section{参考文献}
（同 Markdown 终稿的参考文献列表。）

\section{心得体会}
（随笔，主观感受、批评或建议，内容不限。）

\end{document}
